\section{AMC 10 10A 2010}
        \begin{enumerate}
        	\item (\textit{AMC 10 10A 2010}) Mary’s top book shelf holds five books with the following widths, in centimeters: $6$, $\dfrac{1}{2}$ , $1$, $2.5$, and $10$. What is the average book width, in centimeters?


 $\mathrm{(A)}\ 1 \qquad \mathrm{(B)}\ 2 \qquad \mathrm{(C)}\ 3 \qquad \mathrm{(D)}\ 4 \qquad \mathrm{(E)}\ 5$


 
 \textbf{Answer: }D



\textbf{Solution}

To find the average, we add up the widths $6$, $\dfrac{1}{2}$, $1$, $2.5$, and $10$, to get a total sum of $20$. Since there are $5$ books, the average book width is $\frac{20}{5}=4$ The answer is $\boxed{D}$.

 


	\item (\textit{AMC 10 10A 2010}) Four identical squares and one rectangle are placed together to form one large square as shown. The length of the rectangle is how many times as large as its width? 


 
                    \begin{figure}[H]
                        \centering
                        \includegraphics[width=0.5\textwidth]{Problem/AMC_10_10A_2010_2_Problem_Figure_1.png}
                    \end{figure}
                    

 $\mathrm{(A)}\ \dfrac{5}{4} \qquad \mathrm{(B)}\ \dfrac{4}{3} \qquad \mathrm{(C)}\ \dfrac{3}{2} \qquad \mathrm{(D)}\ 2 \qquad \mathrm{(E)}\ 3$


 
 \textbf{Answer: }B



\textbf{Solution 1}

Let the length of the small square be $x$, intuitively, the length of the big square is $4x$. It can be seen that the width of the rectangle is $3x$. Thus, the length of the rectangle is $4x/3x = 4/3$ times as large as the width. The answer is $\boxed{B}$.

 



\textbf{Solution 2}

We can say the area of one small square is $x^2$, so $\dfrac{1}{4}$ of the area of the large square is $4x^2$ so the area of the large square is $16x^2$, so each side is $4x$ so the length of the rectangle is $4x$ and the width of the rectangle is also $4x$. $4x-x=3x$ so $\dfrac{4x}{3x}=\dfrac{4}{3}$.

 



\textbf{Solution 3}

Let the side length of one of the squares equal $1$. Then, the width of the rectangle will be $4$, and since the width of the rectangle is the same as the length of the entire shape, the length of the rectangle is $4 - 1 = 3$. The ratio between the two is therefore $\frac{4}{3}$ , so our answer is $\boxed{B}$. 

 


	\item (\textit{AMC 10 10A 2010}) Tyrone had $97$ marbles and Eric had $11$ marbles. Tyrone then gave some of his marbles to Eric so that Tyrone ended with twice as many marbles as Eric. How many marbles did Tyrone give to Eric?


 $\mathrm{(A)}\ 3 \qquad \mathrm{(B)}\ 13 \qquad \mathrm{(C)}\ 18 \qquad \mathrm{(D)}\ 25 \qquad \mathrm{(E)}\ 29$


 
 \textbf{Answer: }D



\textbf{Solution 1}

Let $x$ be the number of marbles Tyrone gave to Eric. Then, $97-x = 2\cdot(11+x)$. Solving for $x$ yields $75=3x$ and $x = 25$. The answer is $\boxed{D}$.

 


 



\textbf{Solution 2}

Since the number of balls Tyrone and Eric have a specific ratio when Tyrone gives some of his balls to Eric, we can divide the total amount of balls by their respective ratios. Altogether, they have $97+11$ balls, or $108$ balls, and it does not change when Tyrone gives some of his to Eric, since none are lost in the process. After Tyrone gives some of his balls away, the ratio between the balls is $2:1$, and added together is $108$. The two numbers add up to $3$ and we divide $108$ by $3$, and the outcome is $36$. This is the number of balls Eric has, and so doubling that results in the number of balls Tyrone has, which is $72$. $97-72$ is $25$, and $36-11$ is $25$, thus proving our statement true, and the answer is $\boxed{D=25}$.

 


 



\textbf{Solution 3}

Plug in the answer choices + some common sense. A and B presumably dont work because they are too low, so we start plugging in from C to E. We find $\boxed{D}$ as the answer.

 


	\item (\textit{AMC 10 10A 2010}) A book that is to be recorded onto compact discs takes $412$ minutes to read aloud. Each disc can hold up to $56$ minutes of reading. Assume that the smallest possible number of discs is used and that each disc contains the same length of reading. How many minutes of reading will each disc contain?


 $\mathrm{(A)}\ 50.2 \qquad \mathrm{(B)}\ 51.5 \qquad \mathrm{(C)}\ 52.4 \qquad \mathrm{(D)}\ 53.8 \qquad \mathrm{(E)}\ 55.2$


 
 \textbf{Answer: }B



\textbf{Solution}

Assuming that there are fractions of compact discs, it would take $412/56 ~= 7.357$ CDs to have equal reading time. However, since the number of discs must be a whole number, there are at least 8 CDs, in which case there would be $412/8 = 51.5$ minutes of reading on each of the 8 discs. The answer is $\boxed{B}$.

 



\textbf{Solution 2}

We look at the options, and see which one is divisible by 412. We try A, and see that you don't get a whole number. Next, we try B. We see that the B is divisible by 412 and find that $\boxed{B}$ is our answer.

 -ILoveMath31415926535

 


	\item (\textit{AMC 10 10A 2010}) The area of a circle whose circumference is $24\pi$ is $k\pi$. What is the value of $k$?


 $\mathrm{(A)}\ 6 \qquad \mathrm{(B)}\ 12 \qquad \mathrm{(C)}\ 24 \qquad \mathrm{(D)}\ 36 \qquad \mathrm{(E)}\ 144$


 
 \textbf{Answer: }E



\textbf{Solution}

If the circumference of a circle is $24\pi$, the radius would be $12$. Since the area of a circle is $\pi r^2$, the area is $144\pi$. The answer is $\boxed{E}$.

 


 



\textbf{Solution 2}

By definition, $\pi$ is the ratio of the circumference to the diameter. Since the circumference is $24\pi$, the diameter must be $24$ and the radius is $12$. Therefore, by the area of circle formula $A=\pi r^{2}$ the area is $12^{2}\pi=144\pi$ and $k=144 \Longrightarrow \boxed{\textbf{(E)} 144}$.

 


	\item (\textit{AMC 10 10A 2010}) For positive numbers $x$ and $y$ the operation $x\spadesuit y$ is defined as 


 
 \[x\spadesuit y = x -\dfrac{1}{y}\]


 What is $2\spadesuit (2\spadesuit 2)$?


 $\mathrm{(A)}\ \dfrac{2}{3} \qquad \mathrm{(B)}\ 1 \qquad \mathrm{(C)}\ \dfrac{4}{3} \qquad \mathrm{(D)}\ \dfrac{5}{3} \qquad \mathrm{(E)}\ 2$


 
 \textbf{Answer: }C



\textbf{Solution}

$\spadesuit (2,2) =2-\frac{1}{2} =\frac{3}{2}$. Then, $\spadesuit \left(2,\frac{3}{2}\right)$ is $2-\frac{1}{\frac{3}{2}} = 2- \frac{2}{3} = \frac{4}{3}$The answer is $\boxed{C}$

 



\textbf{Solution 2}


 \[\spadesuit (x, y) \text{is defined as } x - \frac{1}{y} \text{. Hence } \spadesuit (2,\spadesuit(2, 2)) =2 - \frac{1}{\spadesuit (2, 2)} =  2 - \frac{1}{2 - \frac{1}{2}}=2-\frac{1}{\frac{3}{2}}=2-\frac{2}{3}=\frac{4}{3}\Longrightarrow \boxed{\textbf{(C) } \frac{4}{3}}\]

 


	\item (\textit{AMC 10 10A 2010}) Crystal has a running course marked out for her daily run. She begins this run by heading due north for one mile. She then runs northeast for one mile, then southeast for one mile. The last portion of her run takes her on a straight line back to where she started. How far, in miles, is this last portion of her run?


 $\mathrm{(A)}\ 1 \qquad \mathrm{(B)}\ \sqrt{2} \qquad \mathrm{(C)}\ \sqrt{3} \qquad \mathrm{(D)}\ 2 \qquad \mathrm{(E)}\ 2\sqrt{2}$


 
 \textbf{Answer: }C



\textbf{Solution 1}

Crystal first runs north for one mile. Changing directions, she runs northeast for another mile. The angle difference between north and northeast is 45 degrees. She then switches directions to southeast, meaning a 90 degree angle change. The distance now from traveling north for one mile, and her current destination is $\sqrt{2}$ miles, because it is the hypotenuse of a 45-45-90 triangle with side length one (mile). Therefore, Crystal's distance from her starting position, x, is equal to $\sqrt{((\sqrt{2})^2+1^2)}$, which is equal to $\sqrt{3}$. The answer is $\boxed{C}$

 

                    \begin{figure}[H]
                        \centering
                        \includegraphics[width=0.5\textwidth]{Solution/AMC_10_10A_2010_7_Solution_Figure_1.png}
                    \end{figure}
                    

 


 



\textbf{Solution 2}

Crystal runs north one mile, then her next two moves can be broken up into four individual moves: for her northeast section, it forms a $45-45-90$ triangle whose legs are each $\frac{\sqrt{2}}{2}$. For her southeast section, it is also a $45-45-90$ triangle whose legs are each $\frac{\sqrt{2}}{2}$. Notice that the two of the legs cancel each other out; she moves north $\frac{\sqrt{2}}{2}$ units and also south $\frac{\sqrt{2}}{2}$ units. So her net movement these two moves is $\sqrt{2}$ to the right. Finally, after the third move, she is at the corner of a right triangle with legs $1$ and $\sqrt{2}$. Using the Pythagorean theorem, $d^{2}=1^{2}+\left(\sqrt{2}\right)^{2}=1+2=3$ and $d=\sqrt{3} \Longrightarrow \boxed{\textbf{(C) } \sqrt{3}}$.

 


	\item (\textit{AMC 10 10A 2010}) Tony works $2$ hours a day and is paid $$0.50$ per hour for each full year of his age. During a six month period Tony worked $50$ days and earned $$630$. How old was Tony at the end of the six month period?


 $\mathrm{(A)}\ 9 \qquad \mathrm{(B)}\ 11 \qquad \mathrm{(C)}\ 12 \qquad \mathrm{(D)}\ 13 \qquad \mathrm{(E)}\ 14$


 
 \textbf{Answer: }D



\textbf{Solution 1}

Tony works $2$ hours a day and is paid $0.50$ dollars per hour for each full year of his age. This basically says that he gets a dollar for each year of his age. So if he is $12$ years old, he gets $12$ dollars a day. We also know that he worked $50$ days and earned $630$ dollars. If he was $12$ years old at the beginning of his working period, he would have earned $12 * 50 = 600$ dollars. If he was $13$ years old at the beginning of his working period, he would have earned $13 * 50 = 650$ dollars. Because he earned $630$ dollars, we know that he was $13$ for some period of time, but not the whole time, because then the money earned would be greater than or equal to $650$. This is why he was $12$ when he began, but turned $13$ sometime during the six month period and earned $630$ dollars in total. So the answer is $13$.The answer is $\boxed{D}$. We could find out for how long he was $12$ and $13$. $12 \cdot x + 13 \cdot (50-x) = 630$. Then $x$ is $20$ and we know that he was $12$ for $20$ days, and $13$ for $30$ days. Thus, the answer is $13$.

 



\textbf{Solution 2}

Let $x$ equal Tony's age at the end of the period. We know that his age changed during the time period (since $630$ does not evenly divide $50$). Thus, his age at the beginning of the time period is $x - 1$.

 Let $d$ be the number of days Tony worked while his age was $x$. We know that his earnings every day equal his age (since $2 \cdot 0.50 = 1$). Thus, 
 \[x \cdot d + (x - 1)(50 - d) = 630\] 
 \[x\cdot d + 50x - x\cdot d - 50 + d = 630\] 
 \[50x + d = 680\] 
 \[x = \dfrac{680 - d}{50}\]Since $0 < d <50$, $d = 30$. Then we know that $50x = 650$ and $x = \boxed{(D) 13}$

 



\textbf{Solution 3}

Since Tony worked for $50$ days, he has worked for $100$ hours. Let $k$ be his hourly wage. Then, we know that $100k = 630$. 

 Dividing both sides by $50$, we get that $2k$(Daily wage, which is equal to his age) $\approx$ $12.6$. Since we now know that his age was either 12 or 13 during the 6 month span, and we are asked to find his age at the end of this time, the answer is $\boxed{(D) 13}$

 


	\item (\textit{AMC 10 10A 2010}) A \textit{palindrome}, such as $83438$, is a number that remains the same when its digits are reversed. The numbers $x$ and $x+32$ are three-digit and four-digit palindromes, respectively. What is the sum of the digits of $x$?


 $\mathrm{(A)}\ 20 \qquad \mathrm{(B)}\ 21 \qquad \mathrm{(C)}\ 22 \qquad \mathrm{(D)}\ 23 \qquad \mathrm{(E)}\ 24$


 
 \textbf{Answer: }E



\textbf{Solution 1}

$x$ is at most $999$, so $x+32$ is at most $1031$. The minimum value of $x+32$ is $1000$. However, the only palindrome between $1000$ and $1032$ is $1001$, which means that $x+32$ must be $1001$.

 It follows that $x$ is $969$, so the sum of the digits is $\boxed{\textbf{(E)}\ 24}$.

 



\textbf{Solution 2}

For $x+32$ to be a four-digit number, $x$ is in between $968$ and $999$. The palindromes in this range are $969$, $979$, $989$, and $999$, so the sum of the digits of $x$ can be $24$, $25$, $26$, or $27$. Only $\boxed{\textbf{(E)}\ 24}$ is an option, and upon checking, $x+32=1001$ is indeed a palindrome.

 



\textbf{Solution 3}

Since we know $x+32$ to be $1 a a 1$ and the only palindrome that works is $0 = a$, that means $x+32 = 1001$, and so $x = 1001 - 32 = 969$. So, $9$ + $6$ + $9$ = $\boxed{\textbf{(E)}\ 24}$.~songmath20

 


	\item (\textit{AMC 10 10A 2010}) Marvin had a birthday on Tuesday, May 27 in the leap year $2008$. In what year will his birthday next fall on a Saturday?


 $\mathrm{(A)}\ 2011 \qquad \mathrm{(B)}\ 2012 \qquad \mathrm{(C)}\ 2013 \qquad \mathrm{(D)}\ 2015 \qquad \mathrm{(E)}\ 2017$


 
 \textbf{Answer: }E



\textbf{Solution}

There are $365$ days in a non-leap year. There are $7$ days in a week. Since $365 = 52 \cdot 7 + 1$ (or $365\equiv 1 \pmod{ 7}$), the same date (after February) moves "forward" one day in the subsequent year, if that year is not a leap year.

 For example:

 $5/27/08$	Tue

 $5/27/09$	Wed

 $5/27/10$	Thur

 $5/27/11$	Fri

 However, 2012 is a leap year and has $366$ days, and $366 = 52 \cdot 7 + 2$ . This means we move "forward" \textbf{two} days in 2012, and skip past Saturday and straight to Sunday.

 Therefore:$5/27/11$	Fri

 $5/27/12$	Sun

 You can keep counting forward to find that the first time this date falls on a Saturday is in $2017$:

 $5/27/13$	Mon

 $5/27/14$	Tue

 $5/27/15$	Wed

 $5/27/16$	Fri

 $5/27/17$	Sat

 $\boxed{(E) 2017}$

 


	\item (\textit{AMC 10 10A 2010}) The length of the interval of solutions of the inequality $a \le 2x + 3 \le b$ is $10$. What is $b - a$?


 $\mathrm{(A)}\ 6 \qquad \mathrm{(B)}\ 10 \qquad \mathrm{(C)}\ 15 \qquad \mathrm{(D)}\ 20 \qquad \mathrm{(E)}\ 30$


 
 \textbf{Answer: }D



\textbf{Solution 1}

Since we are given the range of the solutions, we must re-write the inequalities so that we have $x$ in terms of $a$ and $b$. 

 $a\le 2x+3\le b$

 Subtract $3$ from all of the quantities:

 $a-3\le 2x\le b-3$

 Divide all of the quantities by $2$.

 $\frac{a-3}{2}\le x\le \frac{b-3}{2}$

 Since we have the range of the solutions, we can make it equal to $10$.

 $\frac{b-3}{2}-\frac{a-3}{2} = 10$

 Multiply both sides by 2.

 $(b-3) - (a-3) = 20$

 Re-write without using parentheses.

 $b-3-a+3 = 20$

 Simplify.

 $b-a = 20$

 We need to find $b - a$ for the problem, so the answer is $\boxed{20\ \textbf{(D)}}$

 



\textbf{Solution 2}

\href{https://artofproblemsolving.com/wiki/index.php?title=Without\_loss\_of\_generality}{Without loss of generality}, let the interval of solutions be $[0, 10]$ (or any real values $[p, 10+p]$). Then, substitute $0$ and $10$ to $x$. This gives $b=23$ and $a=3$. So, the answer is $23-3=\boxed{20\ \textbf{(D)}}$.~ bearjere

 


	\item (\textit{AMC 10 10A 2010}) Logan is constructing a scaled model of his town. The city's water tower stands 40 meters high, and the top portion is a sphere that holds 100,000 liters of water. Logan's miniature water tower holds 0.1 liters. How tall, in meters, should Logan make his tower?


 $\textbf{(A)}\ 0.04 \qquad \textbf{(B)}\ \frac{0.4}{\pi} \qquad \textbf{(C)}\ 0.4 \qquad \textbf{(D)}\ \frac{4}{\pi} \qquad \textbf{(E)}\ 4$


 
 \textbf{Answer: }C



\textbf{Solution}

The water tower holds $\frac{100,000}{0.1} = 1,000,000$ times more water than Logan's miniature. The volume of a sphere is: $V=\dfrac{4}{3}\pi r^3$. Since we are comparing the heights (m), we should compare the radii (m) to find the ratio. Since, the radius is cubed, Logan should make his tower $\sqrt[3]{1,000,000} = 100$ times shorter than the actual tower. This is $\frac{40}{100} = \boxed{0.4}$ meters high, or choice $\textbf{(C)}$.

 Note: The fact that $1\text{ L}=1,000\text{ cm}^3$ doesn't matter since only the ratios are important.

 


	\item (\textit{AMC 10 10A 2010}) Angelina drove at an average rate of $80$ km/h and then stopped $20$ minutes for gas. After the stop, she drove at an average rate of $100$ km/h. Altogether she drove $250$ km in a total trip time of $3$ hours including the stop. Which equation could be used to solve for the time $t$ in hours that she drove before her stop?


 

$\mathrm{(A)}\ 80t + 100\left(\frac{8}{3} -t\right) = 250 \qquad \mathrm{(B)}\ 80t = 250  \qquad \mathrm{(C)}\ 100t = 250$
$\mathrm{(D)}\ 90t = 250 \qquad \mathrm{(E)}\ 80\left(\frac{8}{3} -t\right) + 100t = 250$


 
 \textbf{Answer: }A



\textbf{Solution}

The answer is $A$ because she drove at $80$ kmh for $t$ hours (the amount of time before the stop), and $100$ kmh for $\frac{8}{3}-t$ because she wasn't driving for $20$ minutes, or $\frac{1}{3}$ hours. Multiplying by $t$ gives the total distance, which is $250$ km. Therefore, the answer is $80t+100\left(\frac{8}{3}-t\right)=250$ $\Rightarrow$ $\boxed{(A)}$

 


	\item (\textit{AMC 10 10A 2010}) Triangle $ABC$ has $AB=2 \cdot AC$. Let $D$ and $E$ be on $\overline{AB}$ and $\overline{BC}$, respectively, such that $\angle BAE = \angle ACD$. Let $F$ be the intersection of segments $AE$ and $CD$, and suppose that $\triangle CFE$ is equilateral. What is $\angle ACB$?


 $\textbf{(A)}\ 60^\circ \qquad \textbf{(B)}\ 75^\circ \qquad \textbf{(C)}\ 90^\circ \qquad \textbf{(D)}\ 105^\circ \qquad \textbf{(E)}\ 120^\circ$


 
 \textbf{Answer: }C



\textbf{Solution 1}


                    \begin{figure}[H]
                        \centering
                        \includegraphics[width=0.5\textwidth]{Solution/AMC_10_10A_2010_14_Solution_Figure_1.png}
                    \end{figure}
                    

 
Let $\angle BAE = \angle ACD = x$.

 
 \begin{align*}\angle BCD &= \angle AEC = 60^\circ\\  \angle EAC + \angle FCA + \angle ECF + \angle AEC &= \angle EAC + x + 60^\circ + 60^\circ = 180^\circ\\  \angle EAC &= 60^\circ - x\\  \angle BAC &= \angle EAC + \angle BAE = 60^\circ - x + x = 60^\circ\end{align*}


 Since $\frac{AC}{AB} = \frac{1}{2}$ and the angle between the hypotenuse and the shorter side is $60^\circ$, triangle $ABC$ is a $30-60-90$ triangle, so $\angle BCA = \boxed{90^\circ\,\textbf{(C)}}$.

 



\textbf{Solution 2(Trig and Angle Chasing)}

Let 
 \[AB=2a, AC=a\]Let 
 \[\angle BAE=\angle ACD=x\]Because $\triangle CFE$ is equilateral, we get $\angle FCE=60$, so $\angle ACB=60+x$

 Because $\triangle CFE$ is equilateral, we get $\angle CFE=60$. 

 Angles $AFD$ and $CFE$ are vertical, so $\angle AFD=60$. 

 By triangle $ADF$, we have $\angle ADF=120-x$, and because of line $AB$, we have $\angle BDC=60+x$.

 Because Of line $BC$, we have $\angle AEB=120$, and by line $CD$, we have $\angle DFE=120$. 

 By quadrilateral $BDFE$, we have $\angle ABC=60-x$.

 By the Law of Sines:
 \[\frac{\sin(60-x)}{a}=\frac{\sin(60+x)}{2a}\implies 2\sin(60-x)=\sin(60+x)\]

 By the sine addition formula($\sin(a+b)=\sin(a)\cos(b)+\cos(a)\sin(b)$):
 \[2(\sin(60)\cos(-x)+\cos(60)\sin(-x))=\sin(60)\cos(x)+\cos(60)\sin(x)\]

 Because cosine is an even function, and sine is an odd function, we have 
 \[2\sin(60)\cos(x)-2\cos(60)\sin(x)=\sin(60)\cos(x)+\cos(60)\sin(x) \implies \sin(60)\cos(x)=3\cos(60)\sin(x)\]

 We know that $\sin(60)=\frac{\sqrt{3}}{2}$, and $\cos(60)=\frac{1}{2}$, hence 
 \[\frac{\sqrt{3}}{2}\cos(x)=\frac{3}{2}\sin(x)\implies \tan(x)=\frac{\sqrt{3}}{3}\]

 The only value of $x$ that satisfies $60+x<180$(because $60+x$ is an angle of the triangle) is $x=30^{\circ}$. We seek to find $\angle ACB$, which as we found before is $60+x$, which is $90$. The answer is $90, \text{or} \textbf{(C)}$

 -vsamc

 



\textbf{Solution 3 (Similar Triangles)}

Notice that $\angle AEB=\angle AFC = 120^{\circ}$ and $\angle ACF=\angle BAE$. Hence, triangle AEB is similar to triangle CFA. Since $AB=2AC$, $AE=2CF=2FE$, as triangle CFE is equilateral. Therefore, $AF=FE=FC$, and since $\angle AFC=120^{\circ}$, $x=30$. Thus, the measure of $\angle ACE$ equals to $\angle FCE+\angle ACF=90^{\circ}, \text{or} \textbf{(C)}$ -HarryW

 


 



\textbf{Solution 4}

Notice that $\triangle ADF \sim \triangle CDA$ (by AA Similarity.) Since the corresponding angles of a pair of similar triangles are congruent, we have $\angle DAC = \angle DFA = 60^\circ.$ Since $AB = 2 \cdot AC$ and $\angle DFA = 60^\circ,$ we have that $\triangle ABC$ is congruent by SAS to a $30-60-90$ right triangle, which gives the answer $\boxed{\textbf{(C)}}$.~clever14710owl

 


	\item (\textit{AMC 10 10A 2010}) In a magical swamp there are two species of talking amphibians: toads, whose statements are always true, and frogs, whose statements are always false. Four amphibians, Brian, Chris, LeRoy, and Mike live together in this swamp, and they make the following statements.


 Brian: "Mike and I are different species."


 Chris: "LeRoy is a frog."


 LeRoy: "Chris is a frog."


 Mike: "Of the four of us, at least two are toads."


 How many of these amphibians are frogs?


 $\textbf{(A)}\ 0 \qquad \textbf{(B)}\ 1 \qquad \textbf{(C)}\ 2 \qquad \textbf{(D)}\ 3 \qquad \textbf{(E)}\ 4$


 
 \textbf{Answer: }D



\textbf{Solution 1}

Start with Brian. If he is a toad, he tells the truth, hence Mike is a frog. If Brian is a frog, he lies, hence Mike is a frog, too. Thus Mike must be a frog.

 As Mike is a frog, his statement is false, hence there is at most one toad.

 As there is at most one toad, at least one of Chris and LeRoy is a frog. But then the other one tells the truth, and therefore is a toad.

 Hence we must have one toad and $\boxed{\textbf{(D)}\ 3}$ frogs.

 



\textbf{Solution 2 (logical reasoning like solution 1, but a different train of thought)}

Notice that one of Chris and LeRoy must be a frog: if Chris is a frog, then he lies about LeRoy being a frog. Hence LeRoy is a toad. Alternatively, if Chris is a toad, then he tells the truth about LeRoy being a frog.

 Assume Brian is a toad. Then Mike is a frog, and he lies about at least two being toads. This means that none or one of the amphibians is a toad (the opposite of the statement $n\geq2$ is $n<2$, or $n=0, 1$). However, this is absurd because we assumed Brian is a toad, and we know one of Chris and LeRoy is a toad. So our assumption leads to a contradiction.

 Hence Brian must be a frog, and he and Mike are the same species. Mike is also a frog. One of Chris and LeRoy is a frog. There are $3$ frogs in total $\Longrightarrow \boxed{\textbf{(D) } 3}$.

 



\textbf{Solution 3: Casework/Proof by contradiction}

First, we assume that Brian is a toad. Since his statements are truthful, Mike could not be a toad and must be a frog. Since frogs are liars, Brian must be the only toad in the pond, so Chris and Leroy are both frogs. However, they would then be telling the truth about each other, which only toads do. Hence, this case results in a contraction.Hence, Brian is a frog. He is lying, so Mike must also be a frog. Mike's statement is also false, so Chris and Leroy cannot both be toads. As shown in the previous case, them both being frogs is impossible, so one must be a toad and the other a frog. Therefore, there must be a total of 3 frogs.By-Monkey\_King

 


	\item (\textit{AMC 10 10A 2010}) Nondegenerate ${\triangle ABC}$  has integer side lengths, ${\overline{BD}}$ is an angle bisector, $AD = 3$, and $DC=8$. What is the smallest possible value of the perimeter?


 $\textbf{(A)}\ 30 \qquad \textbf{(B)}\ 33 \qquad \textbf{(C)}\ 35 \qquad \textbf{(D)}\ 36 \qquad \textbf{(E)}\ 37$


 
 \textbf{Answer: }B



\textbf{Solution}

By the \href{https://artofproblemsolving.com/wiki/index.php?title=Angle\_Bisector\_Theorem}{Angle Bisector Theorem}, we know that $\frac{AB}{BC} = \frac{3}{8}$. If we use the lowest possible integer values for $AB$ and $BC$ (the lengths of $AD$ and $DC$, respectively), then $AB + BC = AD + DC = AC$, contradicting the \href{https://artofproblemsolving.com/wiki/index.php?title=Triangle\_Inequality}{Triangle Inequality}. If we use the next lowest values ($AB = 6$ and $BC = 16$), the Triangle Inequality is satisfied. Therefore, our answer is $6 + 16 + 3 + 8 = \boxed{33}$, or choice $\textbf{(B)}$.

 



\textbf{Solution 2(Trick)}

We find that $\frac{AB}{BC}=\frac{3}{8}$ by the \href{https://artofproblemsolving.com/wiki/index.php?title=Angle\_Bisector\_Theorem}{Angle Bisector Theorem} so we let the lengths be $3n$ and $8n$, respectively where $n$ is a positive integer. Also since $AD=3$ and $BC=8$, we notice that the perimeter of the triangle is the sum of these, namely $3n+8n+3+8=11n+11.$ This can be factored into $11(n+1)$ and so the sum must be a multiple of $11$. The only answer choice which is a multiple of $11$ is $\boxed{\textbf{(B)} 33}$. ~mathboy282

 


	\item (\textit{AMC 10 10A 2010}) A solid cube has side length $3$ inches. A $2$-inch by $2$-inch square hole is cut into the center of each face. The edges of each cut are parallel to the edges of the cube, and each hole goes all the way through the cube. What is the volume, in cubic inches, of the remaining solid?


 $\textbf{(A)}\ 7 \qquad \textbf{(B)}\ 8 \qquad \textbf{(C)}\ 10 \qquad \textbf{(D)}\ 12 \qquad \textbf{(E)}\ 15$


 
 \textbf{Answer: }A



\textbf{Solution 1}

Imagine making the cuts one at a time. The first cut removes a box $2\times 2\times 3$. The second cut removes two boxes, each of dimensions $2\times 2\times 0.5$, and the third cut does the same as the second cut, on the last two faces. Hence the total volume of all cuts is $12 + 4 + 4 = 20$.

 Therefore the volume of the rest of the cube is $3^3 - 20 = 27 - 20 = \boxed{7\ \textbf{(A)}}$.

 



\textbf{Solution 2}

We can visualize it has a $3$ by $3$ cube with a hollow $2$ by $2$ center, along with 6 "windows" cut out too. The cube formula gives the volume of the $3$ by $3$ cube as $3^3 = 27$ and $2^3 = 8.$ Additionally, the area of the "windows" is the 2x2 cut along with a thickness of $\frac{(3-2)}{2} = \frac{1}{2}.$ The volume of all 6 windows is $6(\frac{1}{2})(2)(2) = 12.$ This gives us our answer of $27 - 8 - 12 = \boxed{7\ \textbf{(A)}}.$

 



\textbf{Solution 3}

We can use \href{https://artofproblemsolving.com/wiki/index.php?title=Principle\_of\_Inclusion-Exclusion}{Principle of Inclusion-Exclusion} (PIE) to find the final volume of the cube.

 There are 3 "cuts" through the cube that go from one end to the other. Each of these "cuts" has $2 \times 2 \times 3=12$ cubic inches. However, we can not just sum their volumes, asthe central $2\times 2\times 2$ cube is included in each of these three cuts. To get the correct result, we can take the sum of the volumes of the three cuts, and subtract the volume of the central cube twice.

 Hence the total volume of the cuts is $3(2 \times 2 \times 3) - 2(2\times 2\times 2) = 36 - 16 = 20$.

 Therefore the volume of the rest of the cube is $3^3 - 20 = 27 - 20 = \boxed{7\ \textbf{(A)}}$.

 



\textbf{Solution 4}

We can visualize the final figure and see a cubic frame. We can find the volume of the figure by adding up the volumes of the edges and corners.

 Each edge can be seen as a $2\times 0.5\times 0.5$ box, and each corner can be seen as a $0.5\times 0.5\times 0.5$ box.

 $12\cdot{\frac{1}{2}} + 8\cdot{\frac{1}{8}} = 6+1 = \boxed{7\ \textbf{(A)}}$.

 


	\item (\textit{AMC 10 10A 2010}) Bernardo randomly picks 3 distinct numbers from the set $\{1,2,3,4,5,6,7,8,9\}$ and arranges them in descending order to form a 3-digit number. Silvia randomly picks 3 distinct numbers from the set $\{1,2,3,4,5,6,7,8\}$ and also arranges them in descending order to form a 3-digit number. What is the probability that Bernardo's number is larger than Silvia's number?


 $\textbf{(A)}\ \frac{47}{72} \qquad \textbf{(B)}\ \frac{37}{56} \qquad \textbf{(C)}\ \frac{2}{3} \qquad \textbf{(D)}\ \frac{49}{72} \qquad \textbf{(E)}\ \frac{39}{56}$


 
 \textbf{Answer: }B



\textbf{Solution 1}

We can solve this by breaking the problem down into $2$ cases and adding up the probabilities.

 
Case $1$: Bernardo picks $9$.If Bernardo picks a $9$ then it is guaranteed that his number will be larger than Silvia's. The probability that he will pick a $9$ is $\frac{1 \cdot \binom{8}{2}}{\binom{9}{3}} = \frac{\frac{8\cdot7}{2}}{\frac{9\cdot8\cdot7}{3\cdot2\cdot1}}=\frac{1}{3}$.

 
Case $2$: Bernardo does not pick $9$.Since the chance of Bernardo picking $9$ is $\frac{1}{3}$, the probability of not picking $9$ is $\frac{2}{3}$.

 If Bernardo does not pick $9$, then he can pick any number from $1$ to $8$. Since Bernardo is picking from the same set of numbers as Silvia, the probability that Bernardo's number is larger is equal to the probability that Silvia's number is larger.

 Ignoring the $9$ for now, the probability that they will pick the same number is the number of ways to pick Bernardo's 3 numbers divided by the number of ways to pick any 3 numbers.

 We get this probability to be $\frac{3!}{8\cdot{7}\cdot{6}} = \frac{1}{56}$

 The probability of Bernardo's number being greater is
 \[\frac{1-\frac{1}{56}}{2} = \frac{55}{112}\]

 Factoring the fact that Bernardo could've picked a $9$ but didn't:

 
 \[\frac{2}{3}\cdot{\frac{55}{112}} = \frac{55}{168}\]

 Adding up the two cases we get $\frac{1}{3}+\frac{55}{168} = \boxed{\frac{37}{56}\ \textbf{(B)}}$

 


	\item (\textit{AMC 10 10A 2010}) Equiangular hexagon $ABCDEF$ has side lengths $AB=CD=EF=1$ and $BC=DE=FA=r$. The area of $\triangle ACE$ is $70\%$ of the area of the hexagon. What is the sum of all possible values of $r$?


 $\textbf{(A)}\ \frac{4\sqrt{3}}{3} \qquad \textbf{(B)} \frac{10}{3} \qquad \textbf{(C)}\ 4 \qquad \textbf{(D)}\ \frac{17}{4} \qquad \textbf{(E)}\ 6$


 
 \textbf{Answer: }E



\textbf{Solution 1}

It is clear that $\triangle ACE$ is an equilateral triangle. From the \href{https://artofproblemsolving.com/wiki/index.php?title=Law\_of\_Cosines}{Law of Cosines} on $\triangle ABC$, we get that $AC^2 = r^2+1^2-2r\cos{\frac{2\pi}{3}} = r^2+r+1$. Therefore, the area of $\triangle ACE$ is $\frac{\sqrt{3}}{4}(r^2+r+1)$ by area of an equilateral triangle.

 If we extend $BC$, $DE$ and $FA$ so that $FA$ and $BC$ meet at $X$, $BC$ and $DE$ meet at $Y$, and $DE$ and $FA$ meet at $Z$, we find that hexagon $ABCDEF$ is formed by taking equilateral triangle $XYZ$ of side length $r+2$ and removing three equilateral triangles, $ABX$, $CDY$ and $EFZ$, of side length $1$. The area of $ABCDEF$ is therefore

 $\frac{\sqrt{3}}{4}(r+2)^2-\frac{3\sqrt{3}}{4} = \frac{\sqrt{3}}{4}(r^2+4r+1)$.

 
Based on the initial conditions,

 
 \[\frac{\sqrt{3}}{4}(r^2+r+1) = \frac{7}{10}\left(\frac{\sqrt{3}}{4}\right)(r^2+4r+1)\]

 Simplifying this gives us $r^2-6r+1 = 0$. By \href{https://artofproblemsolving.com/wiki/index.php?title=Vieta\%27s\_Formulas}{Vieta's Formulas} we know that the sum of the possible value of $r$ is $\boxed{\textbf{(E)}\ 6}$.

 



\textbf{Solution 2}

Step 1: Use \href{https://artofproblemsolving.com/wiki/index.php?title=Law\_of\_Cosines}{Law of Cosines} in the same manner as the previous solution to get $AC=\sqrt{r^2+r+1}$. 

 Step 2: $\triangle{ABC}$~$\triangle{CDE}$~$\triangle{EFA}$ via SAS congruency. Using the formula $[ABC]=\frac{ab \sin C}{2}= \frac{r \sqrt{3}}{4}$. The area of the hexagon is equal to $[ACE] + 3[ABC]$. We are given that $70\%$ of this area is equal to $[ACE]$; solving for $AC$ in terms of $r$ gives $AC=\sqrt{7r}$.

 Step 3: $\sqrt{7r}=\sqrt{r^2+r+1} \implies r^2-6r+1=0$ and by \href{https://artofproblemsolving.com/wiki/index.php?title=Vieta\%27s\_Formulas}{Vieta's Formulas} , we get $\boxed{\textbf{E}}$.

 Note: To verify that the quadratic $r^2-6r+1$ has two positive roots, we can either solve for the roots directly or note that discriminant is positive, and there are no negative roots (because then $r^2, -6r, 1$ would all be positive).

 



\textbf{Solution 3}

Find the area of the triangle $ACE$ as how it was done in solution 1. Find the sum of the areas of the congruent  triangles $ABC, CDE, EFA$ as how it was done in solution 2. The sum of these areas is the area of the hexagon, hence the areas of the congruent  triangles $ABC, CDE, EFA$ is $30\%$ of the area of the hexagon. Hence $\frac{7}{3}$ times the latter is equal to the triangle $ACE$. Hence $\frac{7}{3}\cdot\frac{3\sqrt{3}}{4}r=\frac{\sqrt{3}}{4}(r^2+r+1)$. We can simplify this to $7r=r^2+r+1\implies r^2-6r+1=0$. By Vieta's, we get the sum of all possible values of $r$ is $-\frac{-6}{1}=6\text{ or } \boxed{\textbf{E}}$.-vsamc

 



\textbf{Solution 4 (no trig)}


                    \begin{figure}[H]
                        \centering
                        \includegraphics[width=0.5\textwidth]{Solution/AMC_10_10A_2010_19_Solution_Figure_1.png}
                    \end{figure}
                    Drawing not to scale.*To find r, we'll form an equation by finding the area of hexagon $ABCDEF$ and $\triangle{ACE}$ separately. First, connect a directly opposite diagonal such as $CF$. Because the hexagon is equiangular, edges $AF$ and $BC$ protrude from points $A$ and $B$ at the same angle but in opposite directions and equal length. This places points $C$ and $F$ at the same "height" relative to segment $AB$, so segment $CF$ is parallel to $AB$ and also $DE$ (opposite sides are parallel since you rotated through three $120^\circ$ angles between each other).

 Next, we'll calculate the area of isosceles trapezoids $ABCF$ and $CDEF$. Drop a perpendicular down from $B$ to $CF$, and call the intersection to $CF$, $P$. Because adjacent angles between parallel sides in a trapezoid sum to $180^\circ$ and $\angle ABC = 120^\circ$, $\angle BCP = 60^\circ$ meaning $\triangle{BPC}$ is a 30-60-90 right triangle. Now we get $BP = \frac{\sqrt{3}}{2}r$ and $CP = \frac{r}{2}$. $CF$ is equal to $AB + 2CP = r+1$. Hence, the area of trapezoid $ABCF$ is $\frac{1}{2} \cdot \frac{\sqrt{3}}{2}r \cdot (r+1+1) = \frac{\sqrt{3}}{4} r(r+2)$. Similarly, on trapezoid $CDEF$ we find $\triangle EQF$ is 30-60-90, and $EQ = \frac{\sqrt{3}}{2}, FQ = \frac{1}{2}$. Then, area of $CDEF$ is $\frac{\sqrt{3}}{4} (2r + 1)$. Thus, the area of the hexagon is the sum of the two areas, $\frac{\sqrt{3}}{4} (r^2 + 4r + 1)$.

 
Next, notice that $\triangle ACE$ is equilateral. And by Pythagorean theorem in $\triangle CQE$, side $CE = \sqrt{r^2 + r + 1}$. Applying area of equilateral triangle formula, area $\triangle ACE = \frac{\sqrt{3}}{4} (r^2 + r + 1)$. Using this area and the area given from 70\% of the area of the hexagon, we get the equation $\frac{7}{10} \cdot \frac{\sqrt{3}}{4} (r^2 + 4r + 1) = \frac{\sqrt{3}}{4} (r^2 + r + 1)$. Simplifying, we get the quadratic $r^2-6r+1=0$, which by Vieta's yields the answer $6, \boxed{\textbf{E}}$.

 \textbf{Proof Triangle ACE is Equilateral.} We know $\triangle{ABC}$, $\triangle{CDE}$, and $\triangle{EFA}$ are congruent by SAS, so the side opposite the 120 degree angle is also the same (since the triangles are congruent). Thus $\triangle{ACE}$ is equilateral.Q.E.D.~mathboy282

 


	\item (\textit{AMC 10 10A 2010}) A fly trapped inside a cubical box with side length $1$ meter decides to relieve its boredom by visiting each corner of the box. It will begin and end in the same corner and visit each of the other corners exactly once. To get from a corner to any other corner, it will either fly or crawl in a straight line. What is the maximum possible length, in meters, of its path?


 $\mathrm{(A)}\ 4+4\sqrt{2} \qquad \mathrm{(B)}\ 2+4\sqrt{2}+2\sqrt{3} \qquad \mathrm{(C)}\ 2+3\sqrt{2}+3\sqrt{3}$

$\mathrm{(D)}\ 4\sqrt{2}+4\sqrt{3} \qquad \mathrm{(E)}\ 3\sqrt{2}+5\sqrt{3}$


 
 \textbf{Answer: }D



\textbf{Solution 1}

The distance of an interior diagonal in this cube is $\sqrt{3}$ and the distance of a diagonal on one of the square faces is $\sqrt{2}$. It is not possible for the fly to travel any interior diagonal twice, as then it would visit a corner more than once. So, the final sum can have at most  $4$ as the coefficient of $\sqrt{3}$. The other $4$ paths taken can be across a diagonal on one of the faces (try to find a such path!), so the maximum distance traveled is $\textbf{(D)}\ 4\sqrt{2}+4\sqrt{3}$.

 



\textbf{Solution 2}

The path of the fly consists of eight line segments, where each line segment goes from one corner to another corner. The distance of each such line segment is $1$, $\sqrt{2}$, or $\sqrt{3}$.

 The only way to obtain a line segment of length $\sqrt{3}$ is to go from one corner of the cube to the opposite corner. Since the fly visits each corner exactly once, it cannot traverse such a line segment twice. Also, the cube has exactly four such diagonals, so the path of the fly can contain at most four segments of length $\sqrt{3}$. Hence, the length of the fly's path can be at most $\boxed{4 \sqrt{3} + 4 \sqrt{2}\text{, or D.}}$. This length can be achieved by taking the path
 \[A \rightarrow G \rightarrow B \rightarrow H \rightarrow C \rightarrow E \rightarrow D \rightarrow F \rightarrow A.\]

 
                    \begin{figure}[H]
                        \centering
                        \includegraphics[width=0.5\textwidth]{Solution/AMC_10_10A_2010_20_Solution_Figure_1.png}
                    \end{figure}
                    

 


	\item (\textit{AMC 10 10A 2010}) The polynomial $x^3 -ax^2 + bx -2010$ has three positive integer roots. What is the smallest possible value of $a$?


 

$\mathrm{(A)}\ 78 \qquad \mathrm{(B)}\ 88 \qquad \mathrm{(C)}\ 98 \qquad \mathrm{(D)}\ 108 \qquad \mathrm{(E)}\ 118$


 
 \textbf{Answer: }A



\textbf{Solution 1 (Alcumus)}

By \href{https://artofproblemsolving.com/wiki/index.php?title=Vieta\%27s\_Formulas}{Vieta's Formulas}, we know that $a$ is the sum of the three roots of the polynomial $x^3-ax^2+bx-2010$. Again Vieta's Formulas tell us that $2010$ is the product of the three integer roots. Also, $2010$ factors into $2\cdot3\cdot5\cdot67$. But, since the polynomial has only three roots, two of the four prime factors must be multiplied so that we are left with three roots. To minimize $a$, $2$ and $3$ should be multiplied, which means $a$ will be $6+5+67=78$ and the answer is $\boxed{\textbf{(A)78}}$.~JimPickens

 


 


 Note:

 If you are feeling unconfident about $78$, you can try to expand the expression $(x-5)(x-6)(x-67)$ which has roots $5$, $6$, and $67$.
 \[(x-5)(x-6)(x-67)=(x^{2}-11x+30)(x-67)\]
 \[=x\left(x^{2}-11x+30\right)-67\left(x^{2}-11x+30\right)=x^{3}-11x^{2}+30x-67x^{2}+737x-2010=x^{3}-78x^{2}+767x-2010\]

 As we can see, $a=78$, and since this is the least answer choice, we can be confident that the right option is $\boxed{\textbf{(A) } 78}$.

 



\textbf{Solution 2}

We can expand $(x+a)(x+b)(x+c)$ as \begin{align*}(x+a)(x+b)(x+c) &= (x^2+ax+bx+ab)(x+c) \\&= x^3+abx+acx+bcx+abx^2+acx^2+bcx^2+abc \\&= x^3+x^2(a+b+c)+x(ab+ac+bc)+abc \end{align*}


 We do not care about $+bx$ in this case, because we are only looking for $a$.  We know that the constant term is $-2010=-(2\cdot 3\cdot 5\cdot 67)$We are trying to minimize a, such that we have $-ax^2$Since we have three positive solutions, we have $(x-a)(x-b)(x-c)$ as our factors.  We have to combine two of the factors of $2\cdot 3\cdot 5\cdot 67$, and then sum up the $3$ resulting factors.  Since we are minimizing, we choose $2$ and $3$ to combine together.  We get $(x-6)(x-5)(x-67)$ which gives us a coefficient of $x^2$ of $-6-5-67=-78$Therefore $-a=-78$ or $a=\boxed{\textbf{(A)}78}$

 



\textbf{Solution 3}

We want the polynomial $x^3-ax^2+bx-2010$ to have POSITIVE integer roots. That means we want to factor it in to the form $(x-a)(x-b)(x-c).$ We therefore want the prime factorization for $2010$. The prime factorization of $2010$ is $2 \cdot 3 \cdot 5 \cdot 67$. We want the smallest difference of the $3$ roots since by \href{https://artofproblemsolving.com/wiki/index.php?title=Vieta\%27s\_Formulas}{Vieta's Formulas}, $a$ is the sum of the $3$ roots.

 We proceed to factorize it in to $(x-5)(x-6)(x-67)$. Therefore, our answer is $5+6+67$ = $\boxed{\textbf{(A)78}}$.

 
Suggestion for the author: The variables $a$ and $b$ are already defined as coefficients of the cubic polynomial. So, consider using different variables when describing how we want to factor this polynomial in the second sentence. Thanks!

 \textbf{Notes} We can check $5 \cdot 6 \cdot 67$ has the smallest possible sum because of the following reasons:

 $1)$: We don't want to multiply $67$ by anything since that would make the sum of the roots too big.

 $2)$: The smallest number should multiply since that would make the numbers optimally small. Therefore, we want $2$ times $3$.

 \href{https://artofproblemsolving.com/wiki/index.php?title=Vieta\%27s\_Formulas}{Vieta's Formulas}

 


	\item (\textit{AMC 10 10A 2010}) Eight points are chosen on a circle, and chords are drawn connecting every pair of points. No three chords intersect in a single point inside the circle. How many triangles with all three vertices in the interior of the circle are created?


 $\mathrm{(A)}\ 28 \qquad \mathrm{(B)}\ 56 \qquad \mathrm{(C)}\ 70 \qquad \mathrm{(D)}\ 84 \qquad \mathrm{(E)}\ 140$


 
 \textbf{Answer: }A



\textbf{Solution}

To choose a chord, we know that two vertices must be chosen. This implies that for three chords to create a triangle and not intersect at a single point, six points need to be chosen. We also know that for any six points we pick, there is only $1$ way to connect the points such that a triangle is formed in the circle's interior (this is because we want no two chords to be nonintersecting ~Williamgolly). Therefore, the answer is ${{8}\choose{6}}$, which is equivalent to $\boxed{\textbf{(A) }28}$.

 


	\item (\textit{AMC 10 10A 2010}) Each of 2010 boxes in a line contains a single red marble, and for $1 \le k \le 2010$, the box in the $k\text{th}$ position also contains $k$ white marbles. Isabella begins at the first box and successively draws a single marble at random from each box, in order. She stops when she first draws a red marble. Let $P(n)$ be the probability that Isabella stops after drawing exactly $n$ marbles. What is the smallest value of $n$ for which $P(n) < \frac{1}{2010}$?


 $\textbf{(A)}\ 45 \qquad \textbf{(B)}\ 63 \qquad \textbf{(C)}\ 64 \qquad \textbf{(D)}\ 201 \qquad \textbf{(E)}\ 1005$


 
 \textbf{Answer: }A



\textbf{Solution 1}

The probability of drawing a white marble from box $k$ is $\frac{k}{k + 1}$, and the probability of drawing a red marble from box $k$ is $\frac{1}{k+1}$.

 To stop after drawing $n$ marbles, we must draw a white marble from boxes $1, 2, \ldots, n-1,$ and draw a red marble from box $n.$ Thus, 
 \[P(n) = \left(\frac{1}{2} \cdot \frac{2}{3} \cdot \frac{3}{4} \cdots \frac {n - 1}{n}\right) \cdot \frac{1}{n +1} = \frac{1}{n (n + 1)}.\]

 So, we must have $\frac{1}{n(n + 1)} < \frac{1}{2010}$ or $n(n+1) > 2010.$

 Since $n(n+1)$ increases as $n$ increases, we can simply test values of $n$; after some trial and error, we get that the minimum value of $n$ is $\boxed{\textbf{(A) }45}$, since $45(46) = 2070$ but $44(45) = 1980.$

 



\textbf{Solution 2(cheap)}

Do the same thing as Solution 1, but when we get to $n(n+1)>2010$ just test all the answer choices in ascending order(from A to E), and stop when one of the answer choices is greater than $2010$. We get $45(46)=2070$, which is greater than $2010$, so we are done. The answer is $\textbf{(A)}$

 -vsamc

 



\textbf{Solution 3(very cheap)}

Do the same thing as Solution 1, but when we get to $n(n+1)>2010$, remember that in our current era, the square root of the year is in the 40s. $n(n+1)$ is close to $n^2$, so only $\textbf{(A)}$ can be correct.

 


	\item (\textit{AMC 10 10A 2010}) The number obtained from the last two nonzero digits of $90!$ is equal to $n$. What is $n$?


 $\textbf{(A)}\ 12 \qquad \textbf{(B)}\ 32 \qquad \textbf{(C)}\ 48 \qquad \textbf{(D)}\ 52 \qquad \textbf{(E)}\ 68$


 
 \textbf{Answer: }A



\textbf{Solution 1}

We will use the fact that for any integer $n$,
 \begin{align*}(5n+1)(5n+2)(5n+3)(5n+4)&=[(5n+4)(5n+1)][(5n+2)(5n+3)]\\ &=(25n^2+25n+4)(25n^2+25n+6)\equiv 4\cdot 6\\ &=24\pmod{25}\equiv -1\pmod{25}.\end{align*}


 First, we find that the number of factors of $10$ in $90!$ is equal to $\left\lfloor \frac{90}5\right\rfloor+\left\lfloor\frac{90}{25}\right\rfloor=18+3=21$. Let $N=\frac{90!}{10^{21}}$. The $n$ we want is therefore the last two digits of $N$, or $N\pmod{100}$. If instead we find $N\pmod{25}$, we know that $N\pmod{100}$, what we are looking for, could be $N\pmod{25}$, $N\pmod{25}+25$, $N\pmod{25}+50$, or $N\pmod{25}+75$. Only one of these numbers will be a multiple of four, and whichever one that is will be the answer, because $N\pmod{100}$ has to be a multiple of 4.

 If we divide $90!$ by $5^{21}$ to create $M$ by taking out all the factors of $5$ in $90!$, we can write $N$ as $\frac M{2^{21}}$ where
 \[M=1\cdot 2\cdot 3\cdot 4\cdot 1\cdot 6\cdot 7\cdot 8\cdot 9\cdot 2\cdots 89\cdot 18,\]where every multiple of 5 is replaced by the number with all its factors of 5 removed. Specifically, every number in the form $5n$ is replaced by $n$, and every number in the form $25n$ is replaced by $n$.

 The number $M$ can be grouped as follows:

 
 \begin{align*}M= &(1\cdot 2\cdot 3\cdot 4)(6\cdot 7\cdot 8\cdot 9)\cdots(86\cdot 87\cdot 88\cdot 89)\\ &\cdot (1\cdot 2\cdot 3\cdot 4)(6\cdot 7\cdot 8\cdot 9)\cdots (16\cdot 17\cdot 18) \\ &\cdot (1\cdot 2\cdot 3).\end{align*}


 Where the first line is composed of the numbers in $90!$ that aren't multiples of five, the second line is the multiples of five \textbf{and not 25} after they have been divided by five, and the third line is multiples of 25 after they have been divided by 25.

 Using the identity at the beginning of the solution, we can reduce $M$ to

 
 \begin{align*}M&\equiv(-1)^{18} \cdot (-1)^3(16\cdot 17\cdot 18) \cdot (1\cdot 2\cdot 3) \\ &= 1\cdot -21\cdot 6\\ &= -1\pmod{25} =24\pmod{25}.\end{align*}


 Using the fact that $2^{10}=1024\equiv -1\pmod{25}$ (or simply the fact that $2^{21}=2097152$ if you have your powers of 2 memorized), we can deduce that $2^{21}\equiv 2\pmod{25}$. Therefore $N=\frac M{2^{21}}\equiv \frac {24}2\pmod{25}=12\pmod{25}$.

 Finally, combining with the fact that $N\equiv 0\pmod 4$ yields $n=\boxed{\textbf{(A)}\ 12}$.

 



\textbf{Solution 2}

Let $P$ be $90!$ after we truncate its zeros. Notice that $90!$ has exactly (floored) $\left\lfloor\frac{90}{5}\right\rfloor + \left\lfloor\frac{90}{25}\right\rfloor = 21$ factors of 5; thus, 
 \[P = 2^{-21}\cdot 5^{-21}\cdot 90!.\] We shall consider $P$ modulo 4 and 25, to determine its residue modulo 100. It is easy to prove that $P$ is divisible by 4 (consider the number of 2s dividing $90!$ minus the number of 5s dividing $90!$), and so we only need to consider $P$ modulo 25.

 Now, notice that for integers $a, n$ we have
 \[(5n + a)(5n - a) \equiv -a^2 \mod 25.\]

 Thus, for integral a: 
 \[(10a + 1)(10a + 2)(10a + 3)(10a + 4)(10a + 6)(10a + 7)(10a + 8)(10a + 9) \equiv (-1)(-4)(-9)(-16) \equiv 576 \equiv 1 \mod 25.\] Using this process, we can essentially remove all the numbers which had not formerly been a multiple of 5 in $90!$ from consideration.

 Now, we consider the remnants of the 5, 10, 15, 20, ..., 90 not yet eliminated. After factoring out a 2 from each of the 9 even numbers in this sequence, the 10, 20, 30, ..., 90 becomes 1, 2, 3, 4, 1, 6, 7, 8, 9, whose product is 1 mod 25. We have accounted for 9 of the 21 2's thus we still need to multiply by $2^{-12}$. Also, the 5, 15, 25, ..., 85 becomes 1, 3, 1, 7, 9, 11, 13, 3, 17. Multiplying out 1, 3, 1, 7, ..., 17 yields 2 modulo 25, and so we are left to compute $2^{-11}$. Also note that $\phi(25)=5^2-5^1=20$. Euler's Theorem states that $2^{\phi(25)} \equiv 1 \mod 25$, thus $2^{-11} \equiv 2^{-11+20} = 2^9 = 512 \equiv 12 \mod 25$. Because 12 is also a multiple of 4, we can utilize the Chinese Remainder Theorem to show that $P \equiv 12 \mod 100$ and so the answer is $\boxed{12}$.

 



\textbf{Solution 3}

We start of by truncating the $0$s off $90!$, just like Solution 2. Since there are $v_5(90!) = \left\lfloor\frac{90}{5}\right\rfloor+\left\lfloor\frac{90}{25}\right\rfloor = 21$ terminating zeroes, we have the number we obtain from truncating the terminating zeroes at the end of $90!$ will be $\frac{90!}{10^{21}}$.

 By \href{https://artofproblemsolving.com/wiki/index.php?title=Chinese\_Remainder\_Theorem}{Chinese Remainder Theorem}, we can divide the mod $100$ into mod $4$ and mod $25$. We know that there are way more than $23$ $2$s in $90!$, so we have

 
 \[\dfrac{90!}{10^{21}} \equiv 0 \text{ } \text{ (mod } 4 \text{)}\]

 Now, notice that $\frac{90!}{5^{21}}$ is basically $90!$ without any factors that are multiples of $5$.

 
 \[1\cdot 2\cdot 3\cdot 4\cdot 6\cdot 7\cdots 23\cdot 24 \text{ } \text{ (mod } 25 \text{)}\]

 We can rewrite the expression as 

 
 \[(1\cdot 2\cdot 3\cdot 4\cdot 6\cdot 7\cdots 12)^2 \text{ } \text{ (mod } 25 \text{)}\]

 We also have 

 
 \[1\cdot 2\cdot 3\cdot 4\cdot 6\cdot 7\cdots 12 \equiv 8\cdot 8\cdot 9\cdot 11\cdot 12 \equiv 7 \text{ } \text{ (mod } 25 \text{)}\]

 Squaring the expression gets us

 
 \[(1\cdot 2\cdot 3\cdot 4\cdot 6\cdot 7\cdots 12)^2 \equiv 49 \equiv -1 \text{ } \text{ (mod } 25 \text{)}\]

 Notice that there are $3$ of this string of numbers multiplied together occurs four times in $\frac{90!}{10^{21}}$, the fourth being only a partial

 $1\cdot 2\cdot 3\cdot 4\cdots 14 \equiv 24 \equiv -1 \text{ } \text{ (mod } 25 \text{)}$

 Note that as in Solution 1, the residue from extracting the fives is$1\cdot 2\cdot 3\cdot 4\cdots 18 \cdot 1 \cdot 2 \cdot 3 \equiv -1 \text{ } \text{ (mod } 25 \text{)}$

 Multiplying these together gives us

 
 \[\frac{90!}{10^{21}} \equiv \frac{\frac{90!}{5^{21}}}{2^{21}} \equiv \frac{(-1)^5}{2^{21}} \equiv \frac{-1}{2^{21}} \text{ } \text{ (mod } 25 \text{)}\]

 By Euler's Theorem, we have $\frac{-1}{2^{21}} \equiv -13^{21} \equiv -13^{\phi{(25)}}\cdot 13 \equiv -13 \equiv 12 \text{ } \text{ (mod } 25 \text{)}$

 Thus, we want to solve $x \equiv 12 \text{ } \text{ (mod } 25 \text{)}$ and $x \equiv 0 \text{ } \text{ (mod } 4 \text{)}$. An obvious solution is $x\equiv 12$ and by Chinese Remainder Theorem, we have $x\equiv 12$ is the only solution. So thus, $\frac{90!}{10^{21}} \equiv \boxed{12} \text{ } \text{ (mod } 100 \text{)}$

 \textbf{Remark (Chinese Remainder Theorem)} Both solution 1 and solution 2 rely on $n \equiv 12 (\bmod { \quad 25})$, $n \equiv 0 (\bmod { \quad 4})$ to get $n (\bmod { \quad 100})$

 By \href{https://artofproblemsolving.com/wiki/index.php/Chinese\_Remainder\_Theorem}{Chinese Remainder Theorem}, the general solution of the system of $2$ linear congruences is:

 \begin{verbatim}
$n \equiv r_1 (\bmod { \quad m_1})$, $n \equiv r_2 (\bmod { \quad m_2})$, $(m_1, m_2) = 1$Find $k_1$ and $k_2$ such that $k_1 m_1 \equiv 1 (\bmod{ \quad m_2})$, $k_2 m_2 \equiv 1 (\bmod{ \quad m_1})$Then $n \equiv k_2 m_2 r_1 + k_1 m_1 r_2 (\bmod{ \quad m_1 m_2})$
\end{verbatim}

In this problem, $n \equiv 12 (\bmod { \quad 25})$, $n \equiv 0 (\bmod { \quad 4})$:

 \begin{verbatim}
$n \equiv 12 (\bmod{ \quad 25})$, $n \equiv 0 (\bmod{ \quad 4})$, $(25, 4) = 1$$1 \cdot 25 \equiv 1 (\bmod{ \quad 4})$, $19 \cdot 4 \equiv 1 (\bmod{ \quad 25})$Then $n \equiv 19 \cdot 4 \cdot 12 + 1 \cdot 25 \cdot 0 \equiv 12 (\bmod{ \quad 100})$
\end{verbatim}




	\item (\textit{AMC 10 10A 2010}) Jim starts with a positive integer $n$ and creates a sequence of numbers. Each successive number is obtained by subtracting the largest possible integer square less than or equal to the current number until zero is reached. For example, if Jim starts with $n = 55$, then his sequence contains $5$ numbers:


 


 \[\begin{array}{ccccc} {}&{}&{}&{}&55\\ 55&-&7^2&=&6\\ 6&-&2^2&=&2\\ 2&-&1^2&=&1\\ 1&-&1^2&=&0\\ \end{array}\]


 Let $N$ be the smallest number for which Jim’s sequence has $8$ numbers. What is the units digit of $N$?


 $\mathrm{(A)}\ 1 \qquad \mathrm{(B)}\ 3 \qquad \mathrm{(C)}\ 5 \qquad \mathrm{(D)}\ 7 \qquad \mathrm{(E)}\ 9$


 
 \textbf{Answer: }B



\textbf{Solution 1}

We can find the answer by working backwards. We begin with $1-1^2=0$ on the bottom row, then the $1$ goes to the right of the equal's sign in the row above. We find the smallest value $x$ for which $x-1^2=1$ and $x>1^2$, which is $x=2$.

 We repeat the same procedure except with $x-1^2=1$ for the next row and $x-1^2=2$ for the row after that. However, at the fourth row, we see that solving $x-1^2=3$ yields $x=4$, in which case it would be incorrect since $1^2=1$ is not the greatest perfect square less than or equal to $x$ . So we make it a $2^2$ and solve $x-2^2=3$. We continue on using this same method where we increase the perfect square until $x$ can be made bigger than it. When we repeat this until we have $8$ rows, we get:

 
 \[\begin{array}{ccccc}{}&{}&{}&{}&7223\\ 7223&-&84^{2}&=&167\\ 167&-&12^{2}&=&23\\ 23&-&4^{2}&=&7\\ 7&-&2^{2}&=&3\\ 3&-&1^{2}&=&2\\ 2&-&1^{2}&=&1\\ 1&-&1^{2}&=&0\\ \end{array}\]

 Hence the solution is the last digit of $7223$, which is $\boxed{\textbf{(B)}\ 3}$.

 Note: We can go up to $167$, and then notice the pattern of units digits alternating between $3$ and $7$, so we do not need to calculate $7223$.

 



\textbf{Solution 2}

Notice that to get the previous term, we must add the smallest square number, (let's call it $n^2$) such that the sum is less than $(n+1)^2$. Otherwise, instead of subtracting $n^2$ from the previous term, we're subtracting a greater square number. 

 
Remember that $(x+1)^2 = x^2 + 2x + 1$. Recall that to find the previous term, we must add a square number such that it is less than the next square number. $a + n^2 < (n+1)^2$. For this to be true, $a < 2n + 1$. What that means is that given a term $a$, we can find the previous term by adding $n^2$ where $n > \frac {a-1}{2}$. 

 
For example, to find the term that precedes $167$, we know that $n>166/2 = 83$. Therefore, $n=84$ and the previous term is $167 + 84^2 = 7223$. The last digit of $7223$ is $3 \Rightarrow \boxed{\textbf{(B)}\ 3}$

 


\end{enumerate}